% ATTEMPT_STATUS: unresolved
% ATTEMPT_ORIGIN: new research, 2026-10-01
% TOP_PROBLEM: {"id": 19, "title": "Sphere Packing in Higher Dimensions", "category_id": 6, "category": {"id": 6, "name": "geometry", "display_name": "Geometry", "description": "Euclidean and non-Euclidean geometry, geometric structures.", "slug": "geometry", "order_index": 6, "created_at": "2026-07-31T15:26:25.671Z"}, "status": "open"}
% FIRST_POSTED: 2026-10-01
% Imported from: unsolvedmath_additional_500.tex; UnsolvedMath ID 19
% !TeX program = lualatex
% Compile twice: lualatex 19.tex
% Self-contained: no external bibliography, images, or \input files.
% Numeric section labels and visible numbers are original UnsolvedMath IDs.
\documentclass[11pt,a4paper]{article}
\usepackage[margin=24mm,headheight=14pt]{geometry}
\usepackage{fontspec}
\setmainfont{Latin Modern Roman}
\setsansfont{Latin Modern Sans}
\setmonofont{Latin Modern Mono}
\usepackage{amsmath,amssymb,mathtools}
\usepackage{unicode-math}
\setmathfont{Latin Modern Math}
\usepackage{microtype}
\usepackage{enumitem,xurl,needspace}
\usepackage[dvipsnames]{xcolor}
\usepackage{hyperref}
\hypersetup{unicode=true,colorlinks=true,linkcolor=MidnightBlue,urlcolor=MidnightBlue,
 pdftitle={UnsolvedMath Problem 19},pdfauthor={Source-annotated compilation},bookmarksnumbered=true}
\usepackage{bookmark,tocloft,titlesec,fancyhdr}
\setlength{\cftsecnumwidth}{4.2em}
\cftsetpnumwidth{2.5em}
\cftsetrmarg{4em}
\titleformat{\section}{\Large\bfseries\raggedright}{\thesection}{0.7em}{}
\pagestyle{fancy}\fancyhf{}
\fancyhead[L]{\small\sffamily UnsolvedMath research attempt}
\fancyhead[R]{\small Original catalogue IDs}
\fancyfoot[C]{\thepage}
\setlength{\parindent}{0pt}
\setlength{\parskip}{5pt plus 1pt}
\setlength{\emergencystretch}{3em}
\setcounter{tocdepth}{1}\setcounter{secnumdepth}{1}
\setlist[itemize]{leftmargin=3em,itemsep=4pt,topsep=4pt,parsep=0pt}
\allowdisplaybreaks
\newcommand{\N}{\mathbb{N}}
\newcommand{\Z}{\mathbb{Z}}
\newcommand{\Q}{\mathbb{Q}}
\newcommand{\R}{\mathbb{R}}
\newcommand{\C}{\mathbb{C}}
\newcommand{\F}{\mathbb{F}}
\newcommand{\A}{\mathbb{A}}
\newcommand{\PP}{\mathbb{P}}
\newcommand{\E}{\mathbb{E}}
\newcommand{\eps}{\varepsilon}
\newcommand{\dd}{\,\mathrm d}
\newcommand{\sref}[2]{\hyperlink{source:#1:#2}{[S#2]}}
\newif\ifshowcatalogue
\showcataloguetrue
\newcommand{\cataloguescope}[1]{\ifshowcatalogue
 \paragraph{Statement recorded in the supplied source.}
 #1\fi}
\begin{document}
% UnsolvedMath ID 19; source catalogue code GEO-002

\setcounter{section}{18}

\section{Sphere Packing in Higher Dimensions}\label{19}

{\small\sffamily UnsolvedMath ID 19 · Geometry}

\subsection{Definitions and mathematical statement}

\paragraph{Shared notation.}Unless a section says otherwise, $\N=\{1,2,\ldots\}$,
$\Z$ is the ring of integers, $\Q,\R,\C$ are the rational, real, and complex
fields, $[N]=\{1,\ldots,\lfloor N\rfloor\}$, and $\log$ is the natural
logarithm. For real functions of a parameter tending to infinity,
$f\ll g$ and $f=O(g)$ mean $|f|\leq Cg$ eventually for some fixed $C$;
$f\gg g$ reverses this comparison; $f=o(g)$ means $f/g\to0$;
$f\sim g$ means $f/g\to1$; and $f\asymp g$ means both order bounds.
A subscript on $O$ or $\ll$ specifies allowed constant dependence.
The expression $n^{o(1)}$ in an upper bound means growth smaller than
$n^\epsilon$ for each fixed $\epsilon>0$. Local definitions govern any
exception, finite-field convention, or use of the same symbol for another
object. An asymptotic claim is not automatically an effective algorithm.



A packing of unit balls in $\R^d$ is a set of centres $X$ with $\|x-y\|_2\geq2$ for distinct $x,y\in X$. Its upper density is
\[\overline\delta(X)=\limsup_{R\to\infty}\frac{\operatorname{vol}_d(B(0,R)\cap\bigcup_{x\in X}B(x,1))}{\operatorname{vol}_d B(0,R)}.\]
Write $\Delta_d=\sup_X\overline\delta(X)$. The target is the exact value of $\Delta_d$ and an optimal packing, without restricting $X$ to a lattice. Scaling the common radius does not change density. \sref{19}{1} \sref{19}{2}

\paragraph{Statement recorded in the supplied source.}
What is the densest packing of congruent spheres in $n$ dimensions for $n \geq 4$? \sref{19}{1}

\paragraph{Scope and interpretation.}

The source asks a family of dimension-by-dimension questions. Its background records dimensions 8 and 24 as solved subcases; those are not presented as new open cases here. \sref{19}{1}

\paragraph{Residual task reported by the archive.}

The embedded review (2026-08-17) records: Resolve the remaining dimensions. \sref{19}{1}

\subsection{Short English statement}

Determine how much of space identical nonoverlapping balls can occupy in each higher dimension, allowing arrangements more general than lattice packings.

\subsection{Research attempt}
\textbf{Status: UNRESOLVED.} We obtain explicit lattice lower bounds and work out a one-parameter Fourier certificate for an upper bound. The certificate can be optimized exactly, but already in dimension four it leaves a substantial gap. These are established elementary constructions and bounds, derived here as an attack on the precise unrestricted-packing target; no improvement of the literature is claimed.

\paragraph{Context and strategy (checked 1 October 2026).}
The primary abstract of Viazovska's paper [S2] confirms optimality of the $E_8$ packing. Cohn--Elkies [E1] supplies the analytic upper-bound theorem used below. We test the simplest Gaussian times a quadratic polynomial, and compare its exact bound with a lattice construction. A numerical candidate without global sign conditions would not suffice.

\paragraph{A lattice construction valid in every dimension.}
Let $v_d=\pi^{d/2}/\Gamma(d/2+1)$ be the volume of the unit ball. The packing $2\mathbb Z^d$ has density $v_d/2^d$. For $d\ge2$ improve this by taking
\[
D_d=\{z\in\mathbb Z^d:\textstyle\sum_jz_j\equiv0\pmod2\}.
\]
The parity homomorphism is surjective, so this sublattice has index and covolume two. A nonzero vector cannot have squared norm one: such a vector has a single coordinate $\pm1$ and odd coordinate sum. The vector $(1,1,0,\ldots,0)$ shows that its least squared norm is two. Consequently $\sqrt2D_d$ has minimum distance two and covolume $2^{d/2+1}$, giving
\[
\Delta_d\ge \frac{v_d}{2^{d/2+1}}.
\]
The density formula follows by averaging the disjoint ball interiors over a fundamental domain; boundary spheres have zero volume. For expanding balls the cells meeting the boundary contribute a relative error tending to zero, because a fundamental domain has bounded diameter. Thus this construction fits the catalogue's ball-window definition of density.

\paragraph{The precise analytic certificate.}
Use $\widehat f(\xi)=\int_{\mathbb R^d}f(x)e^{-2\pi i x\cdot\xi}\,dx$. The Cohn--Elkies theorem implies
\[
\Delta_d\le v_d\frac{f(0)}{\widehat f(0)}
\]
when $f$ is a real radial Schwartz function, $\widehat f\ge0$, $\widehat f(0)>0$, and $f(x)\le0$ for $|x|\ge2$. This statement applies to unrestricted packings; the theorem's passage from periodic packings to the general upper-density definition is an external input.

The periodic calculation explains the certificate. For a union of $M$ distinct translates of a lattice $L$, Poisson summation turns the double sum over pairs of translates into
\[
\frac1{\det L}\sum_{y\in L^*}\widehat f(y)
 \left|\sum_{j=1}^M e^{2\pi i y\cdot x_j}\right|^2.
\]
It is at least $M^2\widehat f(0)/\det L$. On the original side, only the $M$ diagonal terms can be positive, so it is at most $Mf(0)$. Dividing and multiplying by $v_d$ gives the claimed bound for that packing. The nonnegative Fourier transform, not merely decay of $f$, justifies discarding nonzero frequencies.

\paragraph{An explicit certificate and its optimization.}
For $a>0$ set $f_a(x)=(1-|x|^2/4)e^{-\pi a|x|^2}$. Its spatial sign condition is automatic. Differentiating the Gaussian transform, using
$\widehat{|x|^2g}=-(4\pi^2)^{-1}\Delta\widehat g$, gives
\[
\widehat f_a(\xi)=a^{-d/2}e^{-\pi|\xi|^2/a}
 \left(1-\frac{d}{8\pi a}+\frac{|\xi|^2}{4a^2}\right).
\]
Thus $a>d/(8\pi)$ makes it nonnegative and makes its value at zero positive. This yields
\[
\Delta_d\le v_d\frac{a^{d/2}}{1-d/(8\pi a)}.
\]
Writing $b=d/(8\pi)$, the logarithmic derivative of $a^{d/2+1}/(a-b)$ is $(d/2+1)/a-1/(a-b)$. It vanishes precisely at $a=(d+2)/(8\pi)$. The bound diverges both as $a\downarrow b$ and as $a\to\infty$, and the derivative changes sign once, so this is the global minimum within the family. Therefore
\[
\frac{v_d}{2^{d/2+1}}\le\Delta_d
 \le\min\left\{1,\frac{d+2}{2\Gamma(d/2+1)}
             \left(\frac{d+2}{8}\right)^{d/2}\right\}.
\]
The extra upper bound one follows directly from nonoverlap, and is useful if a displayed analytic bound exceeds one.

\paragraph{Adversarial checks and failure of sharpness.}
In dimension four the lower bound is $\pi^2/16$, whereas the optimized upper bound is $27/32$. These are incompatible with a claim of exact determination. At $a=b$, the transform is still nonnegative but its value at zero vanishes, so that endpoint is not an admissible density certificate. The scaling of the lattice and the sign-change radius both use centres at distance two; confusing radius-one and diameter-one conventions would change the density by $2^d$.

There is also a structural obstruction to equality in this family for a lattice: $f_a$ is strictly negative at every distance greater than two, while the lattice has arbitrarily long nonzero vectors. Its spatial Poisson inequality is therefore strict. Adding numerical precision to this one-parameter optimization cannot remove the obstruction. A sharp certificate for a lattice must instead have zeros at its nonzero vector lengths, together with suitable Fourier zeros.

\paragraph{Remaining gap and next step.}
The exact density and an unrestricted optimal packing in the remaining dimensions are unresolved here. The concrete missing step is a matching universal upper certificate (or a denser construction) beyond this Gaussian-quadratic family. A useful continuation is to impose selected lattice-shell zeros in a higher-degree radial polynomial and then prove its two sign conditions on whole half-lines. Checking signs only on a sampled mesh would leave an unjustified step. No finite numerical search or formal proof checking is claimed in this notebook.

\subsection{Sources}

{\small

\begin{itemize}
\item[E1] H. Cohn and N. Elkies, \emph{New upper bounds on sphere packings I}, Annals of Mathematics 157 (2003), 689--714. Primary paper consulted 1 October 2026. \url{https://arxiv.org/abs/math/0110009}.


\item[\hypertarget{source:19:1}{S1}] ULAM.AI, \emph{UnsolvedMath}, supplied \texttt{problems.json}, record 19 (GEO-002), statement and background. The embedded literature-review date is 2026-08-17; that is the archive’s review date, not a fresh certification. Dataset landing page: \url{https://huggingface.co/datasets/ulamai/UnsolvedMath}.

\item[\hypertarget{source:19:2}{S2}] M. Viazovska, The sphere packing problem in dimension 8, Ann. Math. 185 (2017). \url{https://arxiv.org/abs/1603.04246}. Bibliographic information imported from the supplied record.

\end{itemize}

}

\label{end:19}
\end{document}
